\cotefile{157-4}
\scanpage{132}
\mspage{6}
% archive sheet number: 131
(1.16.) Les flèches $i\in C_0$ sont \unsure{co-quarrables}, i.e. pour un diagramme
\diagram{\begin{tikzcd}[ampersand replacement=\&]
A \arrow[r,dash,"u"] \arrow[d,"i"'] \& X \\
B \&
\end{tikzcd}}
avec $i\in C_0$, la somme \unsure{amalgamée} existe.

Mais on peut formuler (1.15) aussi ainsi, indépendamment de l'hypothèse (1.16.) :

(1.1\deleted{\unsure{5|6}}7 \unsure{bis}) Pour \unsure{tt} \unsure{diagramme}
\diagram{\begin{tikzcd}[ampersand replacement=\&]
A \arrow[r,dash,"u"] \arrow[d,"i"'] \& X \arrow[d,"f"] \\
\unsure{B} \arrow[r,"\unsure{v|\nu}"'] \& Y
\end{tikzcd}}
, avec $i\in C_0$, existe \deleted{\unsure{de $f$}} un carré comm. traits pleins
\diagram{\begin{tikzcd}[ampersand replacement=\&]
A \arrow[r,phantom,"\unsure{-|\to}" description,"u"] \arrow[d,"i"'] \& X \arrow[d,"i_1(f)"] \\
B \arrow[r,"v_1"] \arrow[dr,dashed,"v"'] \& X_1(f) \arrow[d,dashed,"p_1(f)"] \\
\& Y
\end{tikzcd}}
tel que $v=p_1(f)v_1$.

Ceci a pour conséquence le
\uline{Corollaire 1} \inserted{(de 1.17)} \inserted{Soit $f\in\mathrm{Fl}\,M$, et} \inserted{$\alpha$ un ordinal $(\in\unsure{\mathcal U|U})$,}
Pour \unsure{tout|tt} \deleted{diagramme}\inserted{carré} commutatif
\diagram{\begin{tikzcd}[ampersand replacement=\&]
A \arrow[r,"\unsure{u_\alpha|u}"] \arrow[d,"i"'] \& X_\alpha(f) \arrow[d,"p_\alpha(f)"] \\
B \arrow[r,"v"'] \& Y
\end{tikzcd}}
il existe un carré commutatif
\diagram{\begin{tikzcd}[ampersand replacement=\&]
A \arrow[r,phantom,"\unsure{\to|-}" description,"u_\alpha"] \arrow[d] \& X_\alpha(f) \arrow[d,"\operatorname{trans}"] \arrow[dd,bend left=45,dashed,"p_\alpha(f)"] \\
B \arrow[r,"v_1"] \arrow[dr,dashed,"v"'] \& X_{\alpha+1}(f) \arrow[d,phantom,"\unsure{\vert|\downarrow}" description,"p_{\alpha+1}(f)"] \\
\& Y
\end{tikzcd}}
tel que
\[v=p_{\alpha+1}(f)v_1.\]
